\chapter{State of the Art and Related Work}\label{ch:related}

\section*{Introduction}
\addcontentsline{toc}{section}{Introduction}

Asking a database a question in plain language is not a new idea, and neither is
wrapping a language model in a loop so it can go and do things. What is a lot
less settled is what happens when you put those two ideas together on top of
operational data that keeps changing while you are querying it, and when the
answers coming out are meant to drive commercial decisions about real customers.
So where does this project actually sit? This chapter places the work against
the research it borrows from, and is upfront about which parts are established
technique and which parts are ours.

\section{From Semantic Parsing to Text-to-SQL}

The systems discussed throughout this chapter rest on the transformer
architecture\cite{Attention} and on the observation that sufficiently large
language models can perform tasks from a handful of in-context examples rather
than task-specific training\cite{gpt3}, later sharpened by tuning models to
follow instructions directly\cite{instructgpt}.

Translating natural language into database queries has been studied for decades,
but the modern framing dates to two benchmark datasets. WikiSQL~\cite{seq2sql}
paired natural-language questions with SQL over single tables, and demonstrated
that sequence-to-sequence models could learn the mapping directly.
Spider~\cite{spider} raised the difficulty substantially by requiring queries
that span multiple joined tables across databases the model has never seen, which
made schema generalisation --- rather than sentence parsing --- the central
problem.

That shift matters a lot for this project. A system that memorises the mapping
from question phrasing to query template becomes useless the moment the schema
changes, and our environment changes continuously. The response in the research
has been to hand the model the schema as context rather than baking it in as
training data.
General-purpose language models proved surprisingly capable at this without
task-specific training~\cite{llm_text2sql_eval}, and systematic benchmarking has
since mapped how prompt design, schema representation and example selection each
move accuracy~\cite{dailsql}.
DIN-SQL~\cite{dinsql} decomposes a hard question into simpler sub-problems and
adds a self-correction pass, reporting that decomposition plus correction
outperforms asking for the whole query in one shot.

Our agent follows that same principle without taking the method wholesale: it
decomposes by taking one action at a time inside a loop rather than asking the
model to plan out sub-queries, and its self-correction is a deterministic check
against the real schema instead of another model call. The reasoning is
something this report comes back to again and again --- a check the model
performs is a check the model can decide to skip.

\section{Language-Model Agents and Tool Use}

Chain-of-thought prompting~\cite{cot} established that eliciting intermediate
reasoning improves performance on multi-step problems, but reasoning alone
cannot consult a database. ReAct~\cite{react} closed that gap by interleaving
reasoning steps with concrete actions against an external environment and
feeding the observations back, which is precisely the cycle at the centre of
Chapter~\ref{ch:agent}. Later work extended the pattern in two directions: searching over several
candidate reasoning paths rather than committing to one~\cite{tot}, and having an
agent critique its own previous attempt before retrying~\cite{reflexion}. Both are
relevant to a loop that must recover from a query returning nothing.
Toolformer~\cite{toolformer} approached the same problem
from the other direction, training a model to decide for itself when to call an
external tool.

We use the ReAct pattern here with one significant practical difference. In the
original version the model writes its action as text and the surrounding program
parses it back out. Our first prototype did exactly that, and hit the predictable
failure: nothing actually forces the model to respect that format, and a small
drift in phrasing breaks the parse silently. Swapping that convention for a
provider-validated tool-calling interface wiped out an entire class of failure
without changing the loop at all --- which is a point about engineering rather
than agent design, but it is the thing that decided whether the system worked.

\section{Retrieval and Grounding}

Retrieval-augmented generation~\cite{rag} established the general remedy for a
model reasoning about knowledge it does not have: retrieve the relevant material
and place it in context rather than relying on parameters. The technique is
normally applied to document collections; here it is applied to schema and query
knowledge, retrieving the subgraph of tables and join paths relevant to a
question along with worked SQL examples.

Retrieval quality rests on established components. BM25~\cite{bm25} provides
lexical matching, which matters when a question names a column almost exactly,
and dense retrieval provides the semantic half~\cite{dpr}, implemented here over
sentence embeddings~\cite{sbert} indexed with FAISS~\cite{faiss} --- which matters
when the question does not name the column at all. Combining both is standard practice, and this system does so.

Where we depart from the retrieval literature is in what happens after
generation. Retrieval lowers the chance of an ungrounded answer, but it does not
detect one, and the tendency of these systems to produce fluent unsupported
content is well documented\cite{hallucination_survey}. The verification layer in Chapter~\ref{ch:agent} --- checking that
every cited number appears in the retrieved rows, and that every claim of absence
has a query behind it --- is a post-hoc deterministic check rather than a
retrieval improvement, and Chapter~\ref{ch:eval} shows both that it catches real
failures and that it misses an entire class of them.

\section{Telecom Analytics and Customer Experience Management}

The operational context is described by the TM~Forum Information
Framework~\cite{tmforum_sid}, which defines the data domains an operator
manages and the OSS/BSS separation the database design in
Chapter~\ref{chapter:BF-DFP} mirrors. Commercial customer-experience platforms,
including the SmartCare product this project was undertaken alongside, operate
over exactly these domains, and their certification against TM~Forum's
architecture~\cite{huawei_tmforum} is what made a faithful synthetic
reconstruction possible without access to proprietary data.

Machine learning applied to churn prediction is long-established, and the
gradient-boosting approach used here~\cite{xgboost} with minority-class
oversampling~\cite{smote} is a conventional treatment of a heavily imbalanced
target. The contribution of this report on that subject is negative and
methodological rather than modelling: Chapter~\ref{ch:eval} documents what
happens when such a model is transferred between subscriber populations without
revalidation.

\section{Positioning}

The decision to have the agent propose rather than act follows established
guidance on human-AI interaction, which recommends making system capability
clear and keeping consequential actions under user control\cite{humanai_guidelines}.

Three families of system are adjacent to this work. Benchmark text-to-SQL
systems optimise query accuracy against static databases with fixed ground
truth. General-purpose business-intelligence assistants translate questions over
a warehouse a customer has configured. And operator-specific analytics platforms
provide deep telecom insight through predefined dashboards rather than open
questioning.

\begin{table}[H]
\centering
\begin{tabular}{lccc}
\toprule
\textbf{Capability} & \textbf{Text-to-SQL} & \textbf{BI assistants} & \textbf{This work} \\
\midrule
Natural-language querying      & Yes & Yes & Yes \\
Cross-domain joins             & Yes & Partial & Yes \\
Data that changes under it     & No  & Yes & Yes \\
Domain-specific grounding      & No  & Partial & Yes \\
Post-hoc answer verification   & No  & No & Yes \\
Refuses under-specified questions & No & No & Yes \\
Human approval before acting   & N/A & N/A & Yes \\
Published accuracy measurement & Yes & Varies & Partial \\
\bottomrule
\end{tabular}
\caption[Capability positioning against adjacent systems]{Capability positioning. The final row is deliberate: benchmark systems
report accuracy on standard datasets, and this work reports a measurement narrow
enough that Chapter~\ref{ch:limits} argues against over-reading it.}
\label{tab:positioning}
\end{table}

The gap we are addressing is the combination rather than any single capability
on that list. Text-to-SQL research has mostly treated the database as static and
the answer as the end of the job. Assistants built over warehouses inherit
whatever grounding the warehouse already has and generally stop at reporting.
The combination we went after --- an operational environment that never stops
moving, a domain the agent is explicitly grounded in, verification applied to
the answer after it has been written, and anything with real consequences routed
to a human --- is what this report contributes, along with an honest account of
where that combination still falls over.
